\documentclass{article}

% ============================================================
% MA 213 In-Class Worksheet Template
% Student-facing worksheet for lecture activities.
% ============================================================

\usepackage[margin=1in]{geometry}
\usepackage{xcolor}
\usepackage{parskip}
\usepackage{array}
\usepackage{enumitem}
\usepackage{listings}
\usepackage{amssymb}

\definecolor{codegray}{gray}{0.97}
\definecolor{huddleblue}{RGB}{214,234,255}
\definecolor{predictionyellow}{RGB}{255,242,200}
\definecolor{debatepink}{RGB}{255,225,225}

\lstset{
  basicstyle=\ttfamily\small,
  columns=fullflexible,
  frame=single,
  framerule=0.4pt,
  backgroundcolor=\color{codegray},
  breaklines=true,
  showstringspaces=false,
  keepspaces=true,
  xleftmargin=0pt,
  xrightmargin=0pt
}

\newcommand{\coursenumber}{MA 213}
\newcommand{\weekplaceholder}{12}
\newcommand{\lectureplaceholder}{31}
\newcommand{\lecturetitle}{Frequentist or Bayesian? Think/Pair/Share}

\newcommand{\nameblank}{\rule{1.75in}{0.4pt}}
\newcommand{\idblank}{\rule{1.55in}{0.4pt}}
\newcommand{\shortblank}{\rule{1.6in}{0.4pt}}
\newcommand{\longblank}{\rule{0.93\linewidth}{0.4pt}}

\newcommand{\responsebox}[2]{
  \noindent\fbox{%
    \begin{minipage}{0.95\linewidth}
      \textbf{#1} #2
      \vspace{0.55in}
    \end{minipage}
  }
}

\newcommand{\linedbox}[1]{
  \noindent\fbox{%
    \begin{minipage}{0.95\linewidth}
      #1
      \vspace{0.18in}

      \longblank

      \vspace{0.18in}
      \longblank

      \vspace{0.18in}
      \longblank
    \end{minipage}
  }
}

\newcommand{\callout}[3]{
  \noindent\colorbox{#1}{%
    \makebox[\dimexpr0.95\linewidth-2\fboxsep\relax][l]{\textbf{#2}}%
  }

  \noindent\fbox{%
    \begin{minipage}{0.93\linewidth}
      #3
    \end{minipage}
  }
}

\newcommand{\classifyrow}[1]{
  \begin{center}
  \renewcommand{\arraystretch}{1.6}
  \begin{tabular}{|c|c|}
  \hline
  \textbf{Frequentist} & \textbf{Bayesian} \\
  \hline
  $\square$ & $\square$ \\
  \hline
  \end{tabular}
  \end{center}
  \textbf{Why?} #1
}

\begin{document}

\begin{center}
  {\LARGE \coursenumber{} Week \weekplaceholder{}: Lecture \lectureplaceholder{}}\\
  {\large \lecturetitle{}}
\end{center}

\begin{enumerate}[label=\arabic*.,leftmargin=*,itemsep=.7em]
  \item \textbf{Your name:} \nameblank \hfill \textbf{BU ID:} \idblank
\end{enumerate}

\textbf{Big idea:} The same real-world situation can often be reasoned about in a Frequentist way \emph{or} a Bayesian way --- they're different philosophies for the same underlying question, not different topics. We will decide, individually and then with a partner, which style of reasoning best describes each scenario below, and why.

\section*{The Scenarios}

For each scenario, decide: is the reasoning being described more naturally \textbf{Frequentist} (long-run frequency of a repeatable procedure; the parameter is fixed) or \textbf{Bayesian} (a degree of belief that gets updated; the parameter is treated as random)?

\begin{enumerate}[leftmargin=*,itemsep=1.1em]

  \item \textbf{An FDA committee approves a drug because its clinical trial produced a $p$-value below 0.05.}
  \classifyrow{}
  \vspace{0.3em}

  \item \textbf{A doctor updates her belief about whether a patient has a disease after each new test result, factoring in how rare the disease is in the general population.}
  \classifyrow{}
  \vspace{0.3em}

  \item \textbf{A pollster reports ``Candidate A leads with a 95\% confidence interval of 48\%--52\%,'' based on the guarantee that this procedure captures the true support 95\% of the time it's used.}
  \classifyrow{}
  \vspace{0.3em}

  \item \textbf{A friend says, ``I'm pretty sure it's going to rain --- the sky looks just like it did the last few times it rained.''}
  \classifyrow{}
  \vspace{0.3em}

  \item \textbf{A factory's quality-control process rejects a batch whenever the defect rate in a sample exceeds a threshold, calibrated so that a \emph{good} batch only gets rejected by mistake 5\% of the time.}
  \classifyrow{}
  \vspace{0.3em}

  \item \textbf{A weather forecaster says, ``There's a 70\% chance of rain tomorrow.''}
  \classifyrow{}

\end{enumerate}

\section*{Step 1: Think (3 mins)}

On your own, work through all six scenarios above and jot down your reasoning.

\section*{Step 2: Pair (5 mins)}

\begin{enumerate}[label=\arabic*.,leftmargin=*,itemsep=.7em, start=2]
  \item \textbf{Partner's name:} \nameblank \hfill \textbf{BU ID:} \idblank
\end{enumerate}

Compare your answers to all six scenarios with your partner. Use these questions to guide your conversation:

\begin{itemize}
  \item Did you and your partner classify any scenario differently? Talk through both readings --- is one of you wrong, or can the scenario genuinely be read either way?
  \item Scenario 6 (the weather forecaster) is a classic hard case. Can you make a genuine argument for \emph{both} the Frequentist and Bayesian reading of ``70\% chance of rain''? What would each one actually mean?
  \item Pick one scenario you both called Frequentist. What would have to be true about the underlying process for that Frequentist interpretation to be meaningful (a repeatable procedure, a long-run frequency)?
\end{itemize}

\textbf{Our thoughts:}
\vfill

\end{document}
